\documentclass[12pt,a4paper]{ctexart}

\usepackage[a4paper,margin=2.2cm]{geometry}
\usepackage{amsmath,amssymb,bm,booktabs,array,longtable}
\usepackage{hyperref}
\usepackage{enumitem}

\hypersetup{
  colorlinks=true,
  linkcolor=blue,
  urlcolor=blue
}

\title{本项目领域自适应数学原理教程}
\author{面向 \texttt{seu\_undergraduate\_final\_year\_project} 仓库}
\date{依据当前工作区代码与配置整理}

\newcommand{\R}{\mathbb{R}}
\newcommand{\E}{\mathbb{E}}
\newcommand{\B}{\mathcal{B}}
\newcommand{\D}{\mathcal{D}}
\newcommand{\Y}{\mathcal{Y}}
\newcommand{\softmax}{\mathrm{softmax}}
\newcommand{\CE}{\mathcal{L}_{\mathrm{ce}}}

\begin{document}

\maketitle
\tableofcontents

\section{写在前面}

这份教程不是泛泛而谈的“领域自适应概论”，而是严格对照当前仓库里的实现来讲数学。它主要依据以下文件整理：
\begin{itemize}[leftmargin=2em]
  \item 方法实现：\texttt{src/methods/*.py}
  \item 损失实现：\texttt{src/losses/domain.py}
  \item 共享骨干：\texttt{src/backbones/fcn.py}
  \item 默认方法配置：\texttt{configs/method/*.yaml}
  \item 当前主线实验配置：\texttt{configs/experiment/*.yaml}
\end{itemize}

因此，下文遇到“标准论文公式”和“仓库当前工程实现”不完全一致时，一律以仓库代码为准。

\section{问题定义与仓库中的数学对象}

\subsection{任务是什么}

本项目处理的是无监督领域自适应故障诊断。设源域训练集为
\begin{equation}
\D_s=\{(x_i^s,y_i^s)\}_{i=1}^{n_s},
\end{equation}
目标域训练集为
\begin{equation}
\D_t=\{x_j^t\}_{j=1}^{n_t},
\end{equation}
其中目标域训练阶段没有标签。

在仓库当前数据设定中，不同 TEP 操作模式就是不同领域：
\begin{equation}
\mathcal{M}=\{\mathrm{mode1},\mathrm{mode2},\mathrm{mode3},\mathrm{mode4},\mathrm{mode5},\mathrm{mode6}\}.
\end{equation}
标签空间固定为 29 类：
\begin{equation}
\Y=\{0,1,\dots,28\}.
\end{equation}

\subsection{样本张量到底长什么样}

根据当前 \texttt{data/benchmark/manifest.json}，6 个域都共享相同单样本长度：
\begin{equation}
x \in \R^{600\times 34}.
\end{equation}
仓库在送入网络前会转为 channels-first，因此真正进入模型的是
\begin{equation}
x' \in \R^{34\times 600}.
\end{equation}

这里 34 是过程变量通道数，600 是时间窗口长度。当前实现对每个域单独做标准化，随后再划分训练/评估索引，这一点由 \nolinkurl{configs/data/te_da.yaml} 与 \nolinkurl{src/datasets/te_torch_dataset.py} 共同决定。

\subsection{统一模型骨架}

无论是 Source Only、CORAL、DAN、DANN、CDAN、DeepJDOT 还是 RCTA，仓库都共享同一个
\begin{equation}
f(x)=h(\phi(x))
\end{equation}
结构：
\begin{itemize}[leftmargin=2em]
  \item 编码器 $\phi$：1D FCN，输出 128 维特征；
  \item 分类头 $h$：隐藏层维度默认 128，输出 29 类 logits。
\end{itemize}

对应到 \texttt{src/backbones/fcn.py}，编码器是三层卷积：
\begin{align}
z^{(1)} &= \mathrm{ReLU}(\mathrm{Norm}(\mathrm{Conv1D}_{9,128}(x'))),\\
z^{(2)} &= \mathrm{ReLU}(\mathrm{Norm}(\mathrm{Conv1D}_{5,256}(z^{(1)}))),\\
z^{(3)} &= \mathrm{ReLU}(\mathrm{Norm}(\mathrm{Conv1D}_{3,128}(z^{(2)}))),
\end{align}
再沿时间维做全局平均池化：
\begin{equation}
f=\phi(x')=\frac{1}{T}\sum_{t=1}^{T} z^{(3)}_{:,t}\in \R^{128}.
\end{equation}
分类头输出 logits：
\begin{equation}
z=h(f)\in \R^{29}, \qquad p=\softmax(z).
\end{equation}

\subsection{统一监督项}

所有方法都保留源域监督分类项。对一个源域小批量 $\B_s$，交叉熵写成
\begin{equation}
\mathcal{L}_{cls}^s
=
-\frac{1}{|\B_s|}
\sum_{(x_i^s,y_i^s)\in \B_s}
\log p_{i,y_i^s}^s.
\end{equation}

需要注意的是，训练器支持多源场景，但当前做法不是“多源专用网络”，而是把多个源 loader 取出的 batch 直接拼起来：
\begin{equation}
\B_s = \bigcup_{m=1}^{M}\B_{s_m}.
\end{equation}
这个行为来自 \texttt{SingleSourceMethodBase.merge\_source\_batches()}。所以从数学上讲，多源只是把监督样本池扩大了；方法本体仍按单一共享特征空间建模。

\subsection{统一的对齐权重调度}

多个方法都会用到对齐权重调度器 \texttt{AdaptationWeightScheduler}。若基础权重为 $\lambda_0$，训练步数为 $t$，最大调度步数为 $T$，则 warm-start 形式为
\begin{equation}
\lambda(t)=\lambda_0\left(\frac{2}{1+\exp(-\alpha\min(t/T,1))}-1\right).
\end{equation}
此外代码还支持 constant 与 linear 两种形式：
\begin{equation}
\lambda_{\mathrm{constant}}(t)=\lambda_0,
\qquad
\lambda_{\mathrm{linear}}(t)=\lambda_0\min(t/T,1).
\end{equation}

\section{六个基线方法}

\subsection{Source Only：只做源域监督}

这是最基础的参照组。它完全忽略目标域训练样本，只最小化
\begin{equation}
\mathcal{L}_{\mathrm{SO}}=\mathcal{L}_{cls}^s.
\end{equation}

它的重要性不在于“强”，而在于回答一个最根本的问题：如果一个领域自适应方法最终还不如 Source Only，那说明它引入的对齐信号并没有真正缓解跨工况分布偏移。

\subsection{CORAL：对齐二阶统计量}

CORAL 的思想是：如果源域与目标域特征协方差接近，那么两域的整体几何形状会更接近。设源域特征矩阵与目标域特征矩阵分别为
\begin{equation}
F_s\in \R^{n_s\times d}, \qquad F_t\in \R^{n_t\times d}.
\end{equation}
其均值为
\begin{equation}
\mu_s=\frac{1}{n_s}\sum_i f_i^s,\qquad
\mu_t=\frac{1}{n_t}\sum_j f_j^t,
\end{equation}
协方差为
\begin{equation}
C(F)=\frac{(F-\bar F)^\top(F-\bar F)}{\max(n-1,1)}.
\end{equation}

当前仓库实现的 CORAL 损失是
\begin{equation}
\mathcal{L}_{coral}
=
\frac{\|C(F_s)-C(F_t)\|_F^2}{4d^2}
\;+\;
\mathbf{1}_{\mathrm{align\_mean}}
\cdot
\frac{\|\mu_s-\mu_t\|_2^2}{d}.
\end{equation}
于是总目标为
\begin{equation}
\mathcal{L}_{\mathrm{CORAL}}
=
\mathcal{L}_{cls}^s+\lambda(t)\mathcal{L}_{coral}.
\end{equation}

\paragraph{如何理解} 协方差控制“云团形状”，均值项控制“云团中心”。仓库默认方法配置 \nolinkurl{configs/method/coral.yaml} 把 \texttt{align\_mean} 设为 false；但某些实验配置会在 \texttt{method\allowbreak\_overrides} 里显式把它打开。

\subsection{DAN：用 MK-MMD 对齐边缘分布}

DAN 通过多核最大均值差异（MK-MMD）逼近
\begin{equation}
P_s(f)\approx P_t(f).
\end{equation}
给定核函数族 $\{k_u\}_{u=1}^{U}$，高斯核常写成
\begin{equation}
k_u(a,b)=\exp\left(-\frac{\|a-b\|_2^2}{2\sigma_u^2}\right).
\end{equation}
MMD 的经典形式是
\begin{equation}
\mathrm{MMD}^2(F_s,F_t)
=
\left\|
\frac{1}{n_s}\sum_i \varphi(f_i^s)
-
\frac{1}{n_t}\sum_j \varphi(f_j^t)
\right\|_{\mathcal H}^2.
\end{equation}

当前仓库把多个高斯核相加，并通过一个 index matrix 构造线性或非线性估计器，因此总损失可概括为
\begin{equation}
\mathcal{L}_{\mathrm{DAN}}
=
\mathcal{L}_{cls}^s+\lambda(t)\mathcal{L}_{mkmmd}.
\end{equation}

\paragraph{如何理解} 如果说 CORAL 是对齐二阶统计量，那么 DAN 是在 RKHS 里对齐整体现象分布。仓库默认配置里 \texttt{linear\_mmd=true}，而主线 benchmark 配置里会把它改为 \texttt{false}，使用更完整的二次型估计。

\subsection{DANN：域对抗学习}

DANN 的核心是让特征既要有分类性，又要让域判别器分不清来自源域还是目标域。设域判别器为 $D(\cdot)$，定义源域域标签为 1，目标域域标签为 0。当前代码里的域对抗损失写成
\begin{equation}
\mathcal{L}_{adv}
=
\frac{1}{2}
\Big(
\mathrm{BCE}(D(\mathrm{GRL}(F_s)),1)
+
\mathrm{BCE}(D(\mathrm{GRL}(F_t)),0)
\Big).
\end{equation}
总目标为
\begin{equation}
\mathcal{L}_{\mathrm{DANN}}
=
\mathcal{L}_{cls}^s+\lambda(t)\mathcal{L}_{adv}.
\end{equation}

这里关键不是公式本身，而是梯度反转层 GRL。它在前向传播时相当于恒等映射，但反向传播时把梯度乘以负系数：
\begin{equation}
\frac{\partial \mathcal{L}}{\partial f}
\leftarrow
-\lambda_{\mathrm{grl}}
\frac{\partial \mathcal{L}}{\partial f}.
\end{equation}
这样域判别器在努力分域，编码器却被迫学出更域不敏感的特征。

\subsection{CDAN：把类别条件也送进对抗器}

CDAN 认为只对齐边缘分布不够，因为多类别任务里更关键的是“同类跨域靠近、异类跨域不要被混在一起”。因此它不是只把特征 $f$ 送进域判别器，而是构造条件表示
\begin{equation}
g = f \otimes p,
\end{equation}
其中 $p=\softmax(z)$ 是类别概率。

当前仓库支持两种条件映射：
\begin{enumerate}[leftmargin=2em]
  \item 确定性多线性映射：
  \begin{equation}
  g=\mathrm{vec}(p f^\top)\in \R^{Kd};
  \end{equation}
  \item 随机映射：
  \begin{equation}
  g=\frac{(fR_f)\odot(pR_p)}{\sqrt{d_r}}.
  \end{equation}
\end{enumerate}

当前实现里，$p$ 在进入条件映射前会被 detach，也就是 stop-gradient。这意味着对抗项主要通过特征分支起作用，而不会把域判别梯度直接回传到 softmax 概率本身。

若记条件域判别器为 $D_c$，则对抗项可写为
\begin{equation}
\mathcal{L}_{cadv}
=
\mathrm{BCE}(D_c(\mathrm{GRL}(g_s)),1)
+
\mathrm{BCE}(D_c(\mathrm{GRL}(g_t)),0).
\end{equation}

CDAN 还支持熵加权：对每个样本给权重
\begin{equation}
w_i \propto 1+\exp(-H(p_i)),
\end{equation}
熵越低，权重越大。

此外，仓库还实现了可选的 MCC 正则。设目标域 logits 为 $z_t$，温度为 $\tau$，则
\begin{equation}
Q=\softmax(z_t/\tau).
\end{equation}
MCC 会构造类别混淆矩阵并惩罚非对角部分，总体可理解为
\begin{equation}
\mathcal{L}_{mcc}
\uparrow
\Longleftrightarrow
\text{目标域类别边界更混乱。}
\end{equation}

因此，CDAN 的总目标是
\begin{equation}
\mathcal{L}_{\mathrm{CDAN}}
=
\mathcal{L}_{cls}^s
\;+\;
\lambda(t)\mathcal{L}_{cadv}
\;+\;
\lambda_{mcc}\mathcal{L}_{mcc}.
\end{equation}

\subsection{DeepJDOT：联合特征距离与类别代价的最优传输}

DeepJDOT 不是靠“一个判别器”来对齐，而是显式建立源域样本和目标域样本之间的运输计划。设 batch 内源域特征为 $\{f_i^s\}_{i=1}^B$，目标域特征为 $\{f_j^t\}_{j=1}^B$，目标域 logits 为 $\{z_j^t\}_{j=1}^B$，当前实现中的代价矩阵为
\begin{equation}
M_{ij}
=
\lambda_{dist}\frac{\|f_i^s-f_j^t\|_2^2}{d}
+
\lambda_{cl}\big(-\log \softmax(z_j^t)_{y_i^s}\big).
\end{equation}

然后求解
\begin{equation}
\mathcal{L}_{ot}
=
\min_{\gamma\in \Pi(a,b)} \langle \gamma, M\rangle,
\end{equation}
其中 $a,b$ 是均匀边缘分布，$\Pi(a,b)$ 表示满足行列约束的运输计划集合。

因此总目标为
\begin{equation}
\mathcal{L}_{\mathrm{DeepJDOT}}
=
\mathcal{L}_{cls}^s+\lambda(t)\mathcal{L}_{ot}.
\end{equation}

\paragraph{如何理解} DANN 和 CDAN 都是在“分域/混域”的意义上对齐；DeepJDOT 则更像“显式匹配谁该和谁靠近”。在这个仓库里，OT 求解器依赖 POT 包；普通 DeepJDOT 方法可以在 Sinkhorn 与 EMD 间切换。

\section{RCTA：可靠性门控的类条件时序自适应}

\subsection{先看整体思路}

RCTA 不是简单把几个损失加起来，而是把四件事串成闭环：
\begin{enumerate}[leftmargin=2em]
  \item 用 EMA 教师在目标域弱增强视图上产生更稳的软预测；
  \item 用可靠性评分决定哪些目标样本值得信；
  \item 用原型记忆把类结构显式拉出来；
  \item 用课程调度控制这些信号什么时候开始发力、发多大力。
\end{enumerate}

更重要的是，当前仓库中的 RCTA 不是单一固定版本。它支持三种基础对齐分支：
\begin{equation}
\texttt{base\_align}\in\{\texttt{cdan},\texttt{dann},\texttt{deepjdot}\}.
\end{equation}
默认方法配置 \texttt{configs/method/rcta.yaml} 使用的是 \texttt{cdan}，而不是论文草稿里曾经写过的固定 DANN 底座；与此同时，一些 experiment config 又会通过 \texttt{method\_overrides} 把它切成 \texttt{dann}。

\subsection{教师网络与弱强双视图}

对目标样本 $x_j^t$，先构造
\begin{equation}
\tilde x_{j,w}^t=\mathcal{A}_w(x_j^t),\qquad
\tilde x_{j,s}^t=\mathcal{A}_s(x_j^t).
\end{equation}
其中弱增强主要是小幅 jitter 和 scaling；强增强还会加入时间遮挡与通道 dropout。也就是说，RCTA 的“Temporal”不是 RNN 意义上的时间建模，而是围绕固定长度工业时序窗口做视图增强和时间掩蔽。

教师网络参数通过 EMA 更新：
\begin{equation}
\theta_{tea}\leftarrow
\mu\theta_{tea}+(1-\mu)\theta_{stu}.
\end{equation}
教师在弱增强视图上输出
\begin{equation}
p_j^{tea}
=
\softmax\left(\frac{z_j^{tea}}{T_{tea}}\right),
\qquad
\hat y_j=\arg\max_k p_{jk}^{tea}.
\end{equation}
这里 $T_{tea}$ 是教师温度，当前默认配置里是 1.4。

\subsection{可靠性评分：四个分量一起投票}

当前代码给每个目标样本构造四个分量：
\begin{align}
q_j &= \max_k p_{jk}^{tea},\\
e_j &= 1-\frac{-\sum_k p_{jk}^{tea}\log p_{jk}^{tea}}{\log K},\\
u_j &= p_{j,\hat y_j}^{stu},\\
s_j &=
\begin{cases}
\dfrac{1+\cos(f_j^{tea},c_{\hat y_j}^{ref})}{2}, & \hat y_j \text{ 对应参考原型已激活},\\[6pt]
0.5, & \text{否则}.
\end{cases}
\end{align}

它们分别对应：
\begin{itemize}[leftmargin=2em]
  \item $q_j$：教师最高置信度；
  \item $e_j$：逆归一化熵，越大越尖锐；
  \item $u_j$：学生在强增强视图上对教师伪标签的响应；
  \item $s_j$：教师特征与参考原型的相似度。
\end{itemize}

若原始权重为 $\omega_1,\omega_2,\omega_3,\omega_4$，代码会先把它们归一化，再得到
\begin{equation}
r_j
=
\mathrm{clip}_{[0,1]}
\left(
\bar\omega_1 q_j+\bar\omega_2 e_j+\bar\omega_3 u_j+\bar\omega_4 s_j
\right).
\end{equation}
若 \texttt{pseudo\_confidence\_power} 不是 1，还会再做幂次放缩：
\begin{equation}
r_j \leftarrow \mathrm{clip}_{[0,1]}(r_j^\gamma).
\end{equation}

\paragraph{为什么这四项要同时存在} 只看最大概率，容易相信“自信但错”的样本；只看熵，容易忽略类结构；只看原型相似度，早期原型又不稳；只看强增强一致性，可能把过拟合的一致性也当成可靠。四个分量一起用，等于从概率、结构和视图一致性三条线交叉验真。

\subsection{门控机制：不是全局阈值，而是按类筛 top 比例}

门控分两步。

\paragraph{第一步：分数下界} 当前步数为 $t$ 时，门控阈值线性调度为
\begin{equation}
\tau(t)=\tau_{start}+(\tau_{end}-\tau_{start})\min(t/T_\tau,1).
\end{equation}

\paragraph{第二步：按类取前若干比例} 接受比例同样做课程调度：
\begin{equation}
\rho(t)=\rho_{start}+(\rho_{end}-\rho_{start})\min(t/T_\rho,1).
\end{equation}
对每个预测类别 $k$，先定义候选集合
\begin{equation}
I_k(t)=\{j\mid \hat y_j=k,\ r_j\ge \tau(t)\}.
\end{equation}
然后只保留该类内部得分最高的
\begin{equation}
\left\lceil \rho(t)\cdot |I_k(t)| \right\rceil
\end{equation}
个样本进入门控集合 $\Omega_t$。

\paragraph{为什么要按类筛} 如果用全局 top 比例，头部类别会吃掉绝大部分伪标签预算。当前 29 类工业诊断任务里，按类预算能明显减少类别失衡对自训练的污染。

\subsection{基础对齐分支与 MCC}

RCTA 并不是一开始就做对齐。若当前步数 $t< t_{align}$，则
\begin{equation}
\mathcal{L}_{align}=0.
\end{equation}
只有到 \texttt{alignment\_start\_step} 之后，对齐分支才开始工作。

另外，若 \texttt{alignment\_use\_reliable\_only=true} 且当前门控集合非空，那么参与对齐的目标域样本不是整批目标样本，而是
\begin{equation}
\{(\tilde x_{j,w}^t,\hat y_j)\mid j\in \Omega_t\}.
\end{equation}
所以 RCTA 的对齐是“可靠样本优先”的，而不是把全部目标样本一股脑送进对齐器。

于是，
\begin{equation}
\mathcal{L}_{align}\in
\{\mathcal{L}_{CDAN-align},\ \mathcal{L}_{DANN-align},\ \mathcal{L}_{DeepJDOT-align}\}.
\end{equation}

如果启用 MCC，代码把它作用在“实际参与对齐的目标 logits”上，而不是永远对整批目标 logits 计算：
\begin{equation}
\mathcal{L}_{mcc}=\mathrm{MCC}(z_{t,\mathrm{align}}).
\end{equation}

\subsection{伪标签损失}

对门控集合中的目标样本，学生在强增强视图上接受伪标签监督。若启用可靠性加权，当前实现是
\begin{equation}
\mathcal{L}_{pl}
=
\frac{\sum_{j\in \Omega_t} r_j\,\mathrm{CE}(z_j^{stu},\hat y_j)}
{\sum_{j\in \Omega_t} r_j+\varepsilon}.
\end{equation}
若关闭可靠性加权，则退化为普通均值交叉熵。

\subsection{原型记忆：仓库里的版本比论文草稿更具体}

对源域当前 batch，每个类别的 batch 原型是
\begin{equation}
\tilde c_k^{s,batch}
=
\mathrm{Norm}\left(
\frac{1}{|\B_{s,k}|}
\sum_{i:y_i^s=k} f_i^s
\right).
\end{equation}

仓库同时维护两个跨步记忆库：
\begin{equation}
\{c_k^s\}_{k=1}^{K},\qquad
\{c_k^t\}_{k=1}^{K}.
\end{equation}
更新方式是动量平均：
\begin{equation}
c_k \leftarrow \mathrm{Norm}\big(\beta c_k +(1-\beta)\tilde c_k\big).
\end{equation}

但当前实现里的参考原型 $c_k^{ref}$ 不是简单三项线性组合，而是下面这个优先级逻辑：
\begin{equation}
c_k^{ref}=
\begin{cases}
\mathrm{Norm}(c_k^s+c_k^t), & c_k^s,c_k^t \text{ 都已激活},\\
c_k^s, & \text{仅源记忆激活},\\
c_k^t, & \text{仅目标记忆激活},\\
\tilde c_k^{s,batch}, & \text{仅当前源 batch 有该类},\\
\text{inactive}, & \text{否则}.
\end{cases}
\end{equation}
也就是说，源记忆优先于当前源 batch；若目标记忆也存在，再与源记忆归一化融合。

基于此，原型吸引项定义为
\begin{equation}
\mathcal{L}_{att}(F,Y)
=
\frac{1}{|F|}
\sum_i \big(1-\cos(f_i,c_{y_i}^{ref})\big).
\end{equation}

仓库里的 RCTA 原型损失不是单一一项，而是
\begin{equation}
\mathcal{L}_{proto}
=
\mathcal{L}_{att}^{src}
\;+\;
\mathcal{L}_{att}^{tgt,gate}
\;+\;
\frac{1}{2}\mathcal{L}_{soft\_tgt}
\;+\;
\lambda_{sep}\mathcal{L}_{sep}.
\end{equation}

其中，软目标原型项会对整批目标弱增强特征做软加权：
\begin{equation}
w_j^{soft}=
\begin{cases}
1, & r_j\ge r_{rel},\\[4pt]
\dfrac{r_j-r_{semi}}{r_{rel}-r_{semi}}, & r_{semi}\le r_j<r_{rel},\\[8pt]
0, & r_j<r_{semi},
\end{cases}
\end{equation}
\begin{equation}
\mathcal{L}_{soft\_tgt}
=
\frac{
\sum_j w_j^{soft}\big(1-\cos(f_j^{t,w},c_{\hat y_j}^{ref})\big)
}{
\sum_j w_j^{soft}+\varepsilon
}.
\end{equation}

原型分离项则惩罚激活类别原型之间过高的相似度：
\begin{equation}
\mathcal{L}_{sep}
=
\mathrm{mean}_{k<l}
\Big[\max(0,\cos(c_k^{ref},c_l^{ref})-m)\Big].
\end{equation}

\paragraph{一个很关键的实现细节} 当前代码中，目标原型记忆库的更新使用的是所有目标 teacher 特征和 argmax 伪标签，而不是只用门控通过的目标样本。所以“显式监督损失”与“原型记忆更新”采用的目标样本范围并不完全相同。

\subsection{一致性、半可靠样本和不可靠样本}

设学生在强增强视图上的概率为 $p_j^{stu}$，教师概率为 $p_j^{tea}$。当前一致性不是 KL 散度，而是逐样本均方差：
\begin{equation}
v_j=\frac{1}{K}\|p_j^{stu}-p_j^{tea}\|_2^2.
\end{equation}

若 \texttt{consistency\_gate\_only=true}，一致性权重就是是否通过门控；否则使用可靠性分数：
\begin{equation}
w_j^{cons}=
\begin{cases}
\mathbf{1}[j\in \Omega_t], & \text{gate-only},\\
r_j^\eta, & \text{否则}.
\end{cases}
\end{equation}
于是主一致性项为
\begin{equation}
\mathcal{L}_{cons}
=
\frac{\sum_j w_j^{cons}v_j}{\sum_j w_j^{cons}+\varepsilon}.
\end{equation}

仓库还把门控样本进一步分成 reliable / semi-reliable / unreliable 三段。对半可靠样本单独再加一项
\begin{equation}
\mathcal{L}_{semi}
=
\frac{\sum_{j\in \mathcal{S}_{semi}} r_j v_j}{\sum_{j\in \mathcal{S}_{semi}} r_j+\varepsilon},
\end{equation}
对不可靠样本加熵惩罚
\begin{equation}
\mathcal{L}_{ue}
=
\frac{1}{|\mathcal{S}_{unrel}|}
\sum_{j\in \mathcal{S}_{unrel}}
\frac{-\sum_k p_{jk}^{stu}\log p_{jk}^{stu}}{\log K}.
\end{equation}

\subsection{课程化总目标}

伪标签、原型、一致性三项都可以带启动步和预热步。统一写成
\begin{equation}
\lambda_{aux}(t)=
\begin{cases}
0, & t<t_0,\\
\lambda \cdot \min\big((t-t_0)/\Delta t,1\big), & t\ge t_0.
\end{cases}
\end{equation}

于是当前仓库中的 RCTA 总目标准确写法是
\begin{align}
\mathcal{L}_{\mathrm{RCTA}}
=\;&
\mathcal{L}_{cls}^s
+
\lambda_{align}(t)\mathcal{L}_{align}
+
\lambda_{mcc}\mathcal{L}_{mcc}
+
\lambda_{pl}(t)\mathcal{L}_{pl}\nonumber\\
&+
\lambda_{proto}(t)\mathcal{L}_{proto}
+
\lambda_{cons}(t)\mathcal{L}_{cons}
+
\lambda_{semi}\mathcal{L}_{semi}
+
\lambda_{ue}\mathcal{L}_{ue}.
\end{align}

这里最后两项 $\lambda_{semi}$ 和 $\lambda_{ue}$ 在当前实现中是固定系数，不走 ramp 调度。

\subsection{一步训练之后还会发生什么}

当前步反向传播和优化器更新完成后，\texttt{after\_optimizer\_step()} 还会做三件事：
\begin{enumerate}[leftmargin=2em]
  \item 用学生参数 EMA 更新教师；
  \item 更新源原型记忆库与目标原型记忆库；
  \item 步数 \texttt{step\_num} 加 1。
\end{enumerate}

这意味着 RCTA 不是“前向算完 loss 就结束”的静态模型，而是一个带状态演化的训练系统。

\section{把七种方法放在一张图里理解}

\begin{center}
\begin{tabular}{p{2.5cm}p{4.5cm}p{6.5cm}}
\toprule
方法 & 核心对象 & 你应该抓住的数学本质 \\
\midrule
Source Only & $\mathcal{L}_{cls}^s$ & 只在源域学判别边界，不处理域偏移。 \\
CORAL & 协方差、均值 & 对齐特征二阶统计量。 \\
DAN & RKHS 均值嵌入 & 用 MK-MMD 对齐边缘分布。 \\
DANN & 域判别器 + GRL & 让特征对类别有用、对域无用。 \\
CDAN & 特征 $\times$ 类别概率 & 不只混域，还把类别条件一起送进对抗器。 \\
DeepJDOT & 运输计划 $\gamma$ & 显式匹配哪一个源样本更应该对应哪一个目标样本。 \\
RCTA & 可靠性、自训练、原型、课程调度 & 在强基线对齐之上，控制谁能参与目标域监督，以及何时参与、参与多深。 \\
\bottomrule
\end{tabular}
\end{center}

\section{结语}

如果把这套项目压缩成一句话，那就是：
\begin{quote}
共享 FCN 骨干负责学工业时序表征，基线方法各自定义“怎么跨域对齐”，而 RCTA 进一步回答“哪些目标样本值得信、怎样把类别结构稳定地注入训练”。
\end{quote}

读懂本项目数学部分时，最值得反复盯住的不是某个单独公式，而是下面这条层次：
\begin{equation}
\text{监督分类}
\rightarrow
\text{全局对齐}
\rightarrow
\text{条件结构}
\rightarrow
\text{可靠伪监督}
\rightarrow
\text{课程化稳定训练}.
\end{equation}

这正是这个仓库从 Source Only 一路走到 RCTA 的主线。

\end{document}
